\chapter*{Appendix B: Detailed Schedule}
\label{appendix:schedule}
\addcontentsline{toc}{chapter}{Appendix B: Detailed Schedule}

This appendix expands the Phase Two roadmap of Section~\ref{sec:future:roadmap} with per-phase outputs, per-phase acceptance criteria, a risk register, and the planned user-evaluation timeline.

\section*{B.1 Deliverables and Acceptance Criteria per Phase}

Table~\ref{tab:appB:phases} restates the ten sequential phases of Phase Two in tabular form for ease of reading. Effort estimates are in weeks. The total budget is approximately fourteen weeks, with phase~P0 dedicated to the data-quality remediation surfaced in Section~\ref{sec:solution:dataset:quality}.

\begin{table}[!htbp]
\centering
\caption{Phase Two phases, outputs, and acceptance criteria.}
\label{tab:appB:phases}
\footnotesize
\renewcommand{\arraystretch}{1.25}
\begin{tabularx}{\linewidth}{p{0.7cm} p{4.6cm} p{1.5cm} X}
\hline
\textbf{Phase} & \textbf{Output} & \textbf{Effort} & \textbf{Acceptance criterion} \\
\hline
P0 & Data-quality remediation (11 observations in Table~\ref{tab:dataset:quality})
   & 1~week
   & Schedule table deduplicated; out-of-bbox stations triaged; null policy declared for sparse payment fields. \\
P1 & Full Medallion warehouse (Silver + Gold) and chart-style helper
   & 2~weeks
   & Build completes under 30~s; 10 dim tables, 3 fact views, 8 aggregates present; first chart passes readability test. \\
P2 & First end-to-end theme (Demand and Crowding)
   & 2~weeks
   & All four demand charts render; entity-detail panel opens on route click; URL state captures active theme and filters. \\
P3 & Spatial and Connectivity theme
   & 1~week
   & Hex map renders at three cell sizes; 4{,}710 stops render without lag; bus-to-metro layer is interactive. \\
P4 & School-Aware Demand theme
   & 1~week
   & School constellation, treemap, and proximity map render; eight universities geocoded and verified. \\
P5 & Financial and Operations plus Network Health themes
   & 2~weeks
   & P\&L quadrant and market treemap interactive; QoS-1/2/3 computed on GPS sample; DiD estimator runs. \\
P6 & Route Detail theme (time-distance diagram, per-stop box plot, pairwise delay)
   & 2~weeks
   & Time-distance diagram supports zoom and trip hover; box plot filters by date and peak window; advisor review held. \\
P7 & Cross-theme polish (entity panel, navigation, Methodology Side-Panel, URL state)
   & 1~week
   & Entity panel works for all six entity types; Methodology Side-Panel renders for every algorithmic chart. \\
P8 & Stories scrollytelling within source themes
   & 1~week
   & Four stories scroll within source themes; each scene deep-links to the corresponding chart state. \\
P9 & Full-dataset ingest and Phase Two acceptance gates
   & 1~week
   & Fifty schools geocoded; full 50-day GPS loaded; acceptance gates passed; two-minute screencast produced. \\
\hline
\end{tabularx}
\end{table}

\section*{B.2 Risk Register}

Table~\ref{tab:appB:risk} records the six risks tracked across Phase Two. Likelihood and impact are each rated Low, Medium, or High.

\begin{table}[!htbp]
\centering
\caption{Phase Two risk register.}
\label{tab:appB:risk}
\footnotesize
\renewcommand{\arraystretch}{1.25}
\begin{tabularx}{\linewidth}{X p{1.8cm} p{1.6cm} X}
\hline
\textbf{Risk} & \textbf{Likelihood} & \textbf{Impact} & \textbf{Mitigation} \\
\hline
Chart-load latency exceeds the interactive budget on the heatmap or hex-grid endpoint. & Medium & High & Pre-materialise Gold aggregate tables; restrict the analytical surface to read from the aggregate layer only. \\
GPS pull is delayed or unavailable beyond P9. & Medium & Medium & Keep the one-day local GPS extract as a fallback; defer the full reliability scorecard until the pull lands. \\
School geocoding completes below the 60-of-64 target. & Medium & Medium & Backfill to district-level approximation and mark each school with a verified-or-stubbed flag. \\
Visualisation primitives fail to render at catalogue scale. & Low & High & Render each primitive on the largest expected entity set during P2; pre-bin or down-sample where needed. \\
Schema mismatch between AFC and GPS expands during Phase Two. & Medium & Medium & Treat the AFC \texttt{route\_no} as canonical; tag routes as GPS-only or AFC-only on the route dimension. \\
Loss of access to operational data during user evaluation. & Low & High & Snapshot the warehouse to a fixed read-only build before the evaluation window opens. \\
\hline
\end{tabularx}
\end{table}

\section*{B.3 Personnel and Time Tracking}

Phase Two is single-author work. Time is tracked per phase against the estimates listed in Table~\ref{tab:appB:phases}. Two advisor reviews are scheduled (at the end of P2 and at the end of P6), at which points the build is checkpointed and any scope adjustments are negotiated before the next set of phases begins.

\section*{B.4 Post-Phase-Two User-Evaluation Timeline}

The user evaluation is conducted in two waves across the three user roles defined in Section~\ref{sec:solution:problem:roles} (passenger, operating company, public-transport authority). Outcome measures are: per-task completion rate, median time-to-answer, and the number of unsolicited findings per participant per week.

\begin{itemize}
    \item \textbf{Wave one (think-aloud sessions).} One ninety-minute session per participant, with the participant attempting a fixed set of scripted tasks while narrating reasoning aloud. Two to three participants per role.
    \item \textbf{Wave two (free exploration).} Two weeks of self-directed use by the same participants, with weekly check-in calls. Participants log unsolicited findings; the analyst logs task-completion telemetry.
\end{itemize}
